\documentclass[10pt,a4paper]{amsart}
\usepackage[margin=19mm]{geometry}
\usepackage{amsmath,amssymb,mathtools}
\usepackage[scr=rsfs,cal=euler]{mathalfa}
\usepackage{xfrac}
\usepackage[unicode,hidelinks,bookmarks=false]{hyperref}
\newcommand{\ie}{i.e.\ }
\newcommand{\cf}{cf.\ }
\newcommand{\resp}{respectively\ }
\newcommand{\Subsection}{\subsection}
\providecommand{\nobreakdash}{\nobreak\hbox{-}}
\setlength{\emergencystretch}{2em}
\multlinegap0pt
\usepackage{bm}
\usepackage{bbm}
\usepackage{dsfont}
\usepackage{xspace}
\usepackage{cancel}
\usepackage{mathtools}
\usepackage{cleveref}
\usepackage{amsthm}
\makeatletter
\@ifundefined{theorem}{\newtheorem{theorem}{Theorem}}{}
\@ifundefined{proposition}{\newtheorem{proposition}[theorem]{Proposition}}{}
\makeatother
\newcommand{\B}{\mathcal{B}}
\newcommand{\Ftwo}{\mathbb F_2}
\newcommand{\G}{\mathcal{G}}
\newcommand{\MH}{\mathrm{MH}}
\newcommand{\M}{\mathcal{M}}
\newcommand{\NN}{\mathbb{N}}
\newcommand{\T}{\mathcal{T}}
\newcommand{\ZZ}{\mathbb{Z}}
\newcommand{\bF}{\boldsymbol{F}}
\newcommand{\bG}{\boldsymbol{G}}
\newcommand{\bH}{\boldsymbol{H}}
\newcommand{\bb}{\mathfrak{b}}
\newcommand{\eps}{\varepsilon}
\newcommand{\one}{\mathbf{1}}
\newcommand{\rk}{\operatorname{rk}}
\title[Mod 2 magnitude cohomology ring of real hyperplane arrangeme]{Mod 2 magnitude cohomology ring of real hyperplane arrangements}
\author{Ye Liu}
\date{Source: arXiv:2608.24869v1 (CC BY 4.0)}
\begin{document}
\maketitle

\noindent\emph{Source and licence.} Mod 2 magnitude cohomology ring of real hyperplane arrangements. Version arXiv:2608.24869v1 (CC BY 4.0), distributed under the Creative Commons Attribution 4.0 International licence, as recorded in the arXiv licence metadata for this exact version and shown on the abstract page https://arxiv.org/abs/2608.24869v1. Full rights notice: licenses/arxiv-2608.24869-CC-BY-4.0.txt.\par\smallskip

\noindent\textit{Editorial scope.} Counted unit: the Theorem whose source label is \texttt{product-formula} in the article (source0.tex lines 1544--1556, source label \texttt{product-formula}), delivered together with the article's own complete original proof of it (source0.tex lines 1558--1587, one \texttt{proof} environment of 30 lines). No mathematics is written, repaired or re-derived: statement and proof are the article's own text line for line. Required context retained uncounted: layercake (lines 463--465); kappa-support (lines 476--486); sorting-move (lines 675--677); Koizumi-basis (lines 914--923); flat-word-monomial (lines 1492--1494); modular-straightening (lines 1502--1508). Every definition, display and label of those lines is retained unchanged. Omitted from this package and not redistributed: 1--462, 466--475, 487--674, 678--913, 924--1491, 1495--1501, 1509--1543, 1557--1557, 1588--1646. Nothing inside a retained excerpt is omitted or altered: every retained body is delivered exactly as the article's lines are written, so no omission receipt of any kind accompanies this package. Citations: the retained excerpts carry no citation command at all, so no bibliography entry of the article is transcribed and no reference is redistributed. Reference adaptation: none was needed and none was made; every reference inside the delivered unit resolves inside it, so no unresolved reference or citation is produced. Printed slips kept literal: no printed word, symbol, hypothesis or number of the counted statement or of its proof was altered. Layout and class: the article's own class is \texttt{amsart} with packages including biblatex, tikz, amssymb, amsthm, xy, xcolor, thmtools, hyperref, cleveref, graphicx, mathrsfs, tikz-3dplot, mathtools, enumitem; the wrapper is the lane's pinned \texttt{amsart} editorial wrapper at 10pt on A4, which restates the theorem-like environments the retained text needs and adds the article's own macro definitions that the retained text needs (\texttt{\textbackslash B}, \texttt{\textbackslash Ftwo}, \texttt{\textbackslash G}, \texttt{\textbackslash M}, \texttt{\textbackslash MH}, \texttt{\textbackslash NN}, \texttt{\textbackslash T}, \texttt{\textbackslash ZZ}, \texttt{\textbackslash bF}, \texttt{\textbackslash bG}, \texttt{\textbackslash bH}, \texttt{\textbackslash bb}, \texttt{\textbackslash eps}, \texttt{\textbackslash one}, \texttt{\textbackslash rk}), with extra wrapper packages (bm, bbm, dsfont, xspace, cancel, mathtools, cleveref). The delivered numbering is the unit's own. Assets: no retained span contains a figure environment, a graphics-inclusion command, a PDF-page inclusion, a table or a TikZ picture; no image file is copied into this package, the package declares no asset dependency and no image file is redistributed with it. Licence: version arXiv:2608.24869v1 of this article is distributed under the Creative Commons Attribution 4.0 International licence, as recorded in the arXiv licence metadata for this exact version and shown on the abstract page https://arxiv.org/abs/2608.24869v1. A full rights notice is delivered at licenses/arxiv-2608.24869-CC-BY-4.0.txt. Characters: every retained excerpt is pure ASCII, so the delivered engine, XeLaTeX with its default Latin Modern text font, reports no missing character.
\begin{equation}\label{layercake}
 w=\sum_{j=1}^{m(w)}\one_{U_j(w)}.
\end{equation}

\begin{proposition}\label{kappa-support}
For every $w\in\NN^E$,
\[
\kappa_M(w)=\max_{\bb\in\B(M)}w(\bb).
\]
Consequently $\kappa_M$ is subadditive, that is, for $\alpha,\beta\in\NN^E$,
\[
\kappa_M(\alpha+\beta)\leq\kappa_M(\alpha)+\kappa_M(\beta),
\]
where equality holds if and only if $M$ has a basis which simultaneously maximizes $\alpha$ and $\beta$.
\end{proposition}

\begin{equation}\label{sorting-move}
(S_i,S_{i+1})\longmapsto(S_i\cup S_{i+1},S_i\cap S_{i+1}).
\end{equation}

\begin{theorem}\label{Koizumi-basis}
For every $\alpha\in\NN^E$ and $B\in\T$,
\[
\MH_{k,\alpha}(\G;\ZZ)_{\bullet\to B}\cong\begin{cases}
 \ZZ,&k=\kappa_M(\alpha)\text{ and }(\alpha,B)\text{ is tope-admissible},\\
 0,&\text{otherwise}.
 \end{cases}
\]
In particular, all integral magnitude homology groups of $\G$ are torsion-free.
\end{theorem}

\begin{equation}\label{flat-word-monomial}
x_{\bF,C}=x_{F_1,C_1}\smile\cdots \smile x_{F_m,C_m}\in\MH^{d(\bF),w(\bF)}(\G(\M);\Ftwo)_{C_0\to C_m}.
\end{equation}

\begin{proposition}\label{modular-straightening}
If $(\bF,C)$ is clean, then every factor in \eqref{flat-word-monomial} is defined and
\[
x_{\bF,C}=u_{w(\bF),C}.
\]
More precisely, every elementary sorting move \eqref{sorting-move} on $\bF$ is realized by the local modular relation.
\end{proposition}

\begin{theorem}\label{product-formula}
Let $(\alpha,B)$ and $(\beta,C)$ be tope-admissible in $\M$. Then
\begin{equation}\label{product-eq}
u_{\alpha,B}\smile u_{\beta,C}=\begin{cases}
 u_{\alpha+\beta,C}, &\begin{array}{l}
       B=\eps_\beta C,\\
       (\alpha+\beta,C)\text{ is tope-admissible},\\
       \kappa_M(\alpha+\beta)=\kappa_M(\alpha)+\kappa_M(\beta),
 \end{array}\\[6mm]
 0,&\text{otherwise}.
 \end{cases}
\end{equation}
\end{theorem}

\begin{proof}
If $B\neq\eps_\beta C$, the endpoint summands do not compose, so the cup product is zero. Assume $B=\eps_\beta C$. The product has crossing vector $\alpha+\beta$ and cohomological degree $\kappa_M(\alpha)+\kappa_M(\beta)$. If $(\alpha+\beta,C)$ is not tope-admissible, then every cohomology group in the $(\alpha+\beta)$-crossing endpoint summand ending at $C$ is zero by \Cref{Koizumi-basis}. By \Cref{kappa-support}, if $\kappa_M(\alpha+\beta)<\kappa_M(\alpha)+\kappa_M(\beta)$, then \Cref{Koizumi-basis} shows that the target crossing summand vanishes in the cohomological degree of the product. These are the otherwise cases.

Next suppose the three displayed conditions are satisfied. Put $p=\max_{e\in E}\alpha_e$ and $q=\max_{e\in E}\beta_e$ and form flat words
\[
\bF=(F_1,\ldots,F_p),\quad F_i=U_i(\alpha), \quad \text{and} \quad \bG=(G_1,\ldots,G_q),\quad G_j=U_j(\beta).
\]
Also name sign vectors $A_p=B, A_{i-1}=\eps_{F_i}A_i$ and $D_q=C, D_{j-1}=\eps_{G_j}D_j$. By \Cref{layercake}, we have $\alpha=w(\bF)$ and $\beta=w(\bG)$. Moreover $(\bF,B)$ and $(\bG,C)$ are clean since $(\alpha,B),(\beta,C)$ are supposed to be tope-admissible and 
\begin{equation}\label{degree-kappa-eq}
d(\bF)=\sum_{i=1}^p\rk U_i(\alpha)=\kappa_M(\alpha),\quad d(\bG)=\sum_{j=1}^q\rk U_j(\beta)=\kappa_M(\beta),  
\end{equation}
Therefore
\begin{equation}\label{two-factors}
u_{\alpha,B}=x_{\bF,B}=x_{F_1,A_1}\smile\cdots\smile x_{F_p,A_p}, \quad u_{\beta,C}=x_{\bG,C}=x_{G_1,D_1}\smile\cdots\smile x_{G_q,D_q},
\end{equation}
are defined by \Cref{modular-straightening}. Since $D_0=\eps_\beta C=B=A_p$, these two monomials are composable.

Let $\bH=\bF\star\bG=(F_1,\ldots,F_p,G_1,\ldots,G_q)$ be the concatenation word. We claim $(\bH,C)$ is clean. In fact, $w(\bH)=\sum_{i=1}^p\one_{F_i}+\sum_{j=1}^q\one_{G_j}=\alpha+\beta$ and we already assumed $(\alpha+\beta,C)$ is tope-admissible; on the other hand
\[
d(\bH)=\sum_{i=1}^p\rk F_i+\sum_{j=1}^q\rk G_j=\kappa_M(\alpha)+\kappa_M(\beta)=\kappa_M(\alpha+\beta),
\]
where the second equality is \eqref{degree-kappa-eq} and the third is our assumption. Then \Cref{modular-straightening} again shows that
\[
x_{\bH,C}=x_{F_1,A_1}\smile\cdots\smile x_{F_p,A_p}\smile x_{G_1,D_1}\smile\cdots\smile x_{G_q,D_q}
\]
is defined and $u_{\alpha+\beta,C}=x_{\bH,C}$. Then combining with \Cref{two-factors}, we conclude
\[
u_{\alpha+\beta,C}=u_{\alpha,B}\smile u_{\beta,C}.
\]
\end{proof}

\end{document}
